
\documentclass{article}

\usepackage{amsmath}
\usepackage{amssymb}

\title{Mathematical Typesetting in LaTeX}
\author{Siddharth Kumar}
\date{\today}

\begin{document}

\maketitle

\section{Inline and Display Mathematics}

Inline mathematics appears within a sentence, for example
\(x^2 + y^2 = z^2\).

Display mathematics appears on a separate line:

\[
x^2 + y^2 = z^2
\]

\section{Superscripts and Subscripts}

\[
a_i^2 + b_j^3
\]

\[
x_{ij}, \quad x^{n+1}, \quad a_{i,j}^{2}
\]

\section{Greek Letters and Mathematical Symbols}

Greek letters:

\[
\alpha,\ \beta,\ \gamma,\ \delta,\ \lambda,\ \pi,\ \sigma
\]

Relations and set symbols:

\[
a \leq b,\quad x \neq y,\quad
A \subseteq B,\quad A \cup B,\quad A \cap B
\]

Binary operators:

\[
a \times b,\quad A \otimes B,\quad A \oplus B
\]

\section{Fractions and Roots}

\[
\frac{a+b}{c}, \qquad
\sqrt{x}, \qquad
\sqrt[3]{y}
\]

\section{Mathematical Fonts}

Different mathematical fonts can be used to represent expressions.

\[
\mathcal{ABC} \qquad
\mathfrak{ABC} \qquad
\mathrm{ABC}
\]

\[
\mathnormal{ABC} \qquad
\mathbf{ABC} \qquad
\mathsf{ABC}
\]

\section{Mathematical Display Styles}

The same expression can appear in different mathematical styles.

\[
\sum_{i=1}^{n} i
\]

\[
\displaystyle \sum_{i=1}^{n} i
\]

\[
\textstyle \sum_{i=1}^{n} i
\]

\[
\scriptstyle \sum_{i=1}^{n} i
\]

\[
\scriptscriptstyle \sum_{i=1}^{n} i
\]

\section{Summation, Product and Integration}

Summation:

\[
\sum_{i=1}^{n} i = \frac{n(n+1)}{2}
\]

Product:

\[
\prod_{k=1}^{m} k
\]

Definite integral:

\[
\int_{0}^{1} x^2 \, dx = \frac{1}{3}
\]

Other integral symbols:

\[
\oint_C f(z)\,dz
\]


\section{Numbered Equations}

The equation environment automatically numbers equations.

\begin{equation}
    E = mc^2
\end{equation}

\begin{equation}
    \int_0^1 x^2 \, dx = \frac{1}{3}
\end{equation}

\section{Aligned Equations}

The align environment aligns multiple equations
at a common point.

\begin{align}
    a + b &= c \\
    x + y &= 10 \\
    2x + 3y &= 15
\end{align}

\section{Arrays}

The array environment creates columns with
different alignment options.

\[
\begin{array}{lcr}
\text{Left} & \text{Center} & \text{Right} \\
a & b & c \\
1 & 2 & 3
\end{array}
\]

\section{Matrices}

\subsection{Basic Matrix}

\[
\begin{matrix}
1 & 2 \\
3 & 4
\end{matrix}
\]

\subsection{Matrix with Parentheses}

\[
\begin{pmatrix}
1 & 0 \\
0 & 1
\end{pmatrix}
\]

\subsection{Matrix with Square Brackets}

\[
\begin{bmatrix}
a & b \\
c & d
\end{bmatrix}
\]

\subsection{Matrix with Curly Braces}

\[
\begin{Bmatrix}
x & y \\
z & w
\end{Bmatrix}
\]

\subsection{Determinant}

\[
\begin{vmatrix}
1 & 2 \\
3 & 4
\end{vmatrix}
\]

\subsection{Double Vertical Bars}

\[
\begin{Vmatrix}
a & b \\
c & d
\end{Vmatrix}
\]

\section{Description Lists}

The description environment creates lists
containing terms and their explanations.

\begin{description}
    \item[Matrix] A rectangular arrangement of numbers
    in rows and columns.

    \item[Determinant] A scalar value associated
    with a square matrix.

    \item[Integral] A mathematical operation used
    to calculate areas and accumulated quantities.
\end{description}


\end{document}
